\documentclass[conference]{IEEEtran}
\usepackage{amsmath,amssymb,amsfonts}
\usepackage{booktabs}
\usepackage{graphicx}
\usepackage{textcomp}
\usepackage{xcolor}
\usepackage{hyperref}

\begin{document}

\title{Open-Loop L2 is Blind: Closed-Loop NeuroNCAP Reveals\\
That INT4 Quantization Granularity Governs Autonomous Driving Safety}

\author{
\IEEEauthorblockN{Justin Williams, Kishor Datta Gupta, Roy George, Mrinmoy Sarkar}
\IEEEauthorblockA{Clark Atlanta University \quad \texttt{justinwilliamstech693@gmail.com}}
}

\maketitle

\begin{abstract}
Weight-only INT4 quantization of autonomous driving VLAs is commonly reported as ``lossless''
based on open-loop nuScenes L2 trajectory metrics. We show this conclusion is an artifact: the
open-loop nuScenes planning metric is dominated by ego-motion extrapolation (zeroing the ego-state
and history text raises L2 from 0.33\,m to 6.38\,m, $19\times$, while leaving outputs well-formed),
so it cannot detect quantization damage. Moving to the \textbf{NeuroNCAP closed-loop safety
benchmark}---where the model actually drives and perception is on the critical path---reveals that
\emph{naive per-channel} INT4 collapses trajectory generation coherence to 3.7\% (mean across 16
scenario instances, 10 seeds each), while \emph{group-wise} INT4 (group size 128, $\sim$0.1 extra
bits/weight) restores FP16-level coherence (73.5\% vs.\ 76.0\%). The fix is one parameter:
per-channel$\to$group-128 fully recovers generation robustness on the exact perturbations where
per-channel collapses. We report this result on 16 NeuroNCAP scenario instances (14 scenes, 10
seeds each), using real neural-rendered cameras (NeuRAD), with open-loop L2 confirming W4=W8=FP16
at 0.33\,m throughout. The unconfounded metric is \emph{generation coherence} (well-formed
trajectory output rate), not collision rate, which is a confounded proxy. W8 and W4-group128 are
deployment-safe; naive W4 per-channel is not, and open-loop L2 cannot tell you which is which.
\end{abstract}

\begin{IEEEkeywords}
autonomous driving, quantization, NeuroNCAP, closed-loop evaluation, nuScenes, INT4,
group quantization, VLA, safety
\end{IEEEkeywords}

\section{Introduction}
\label{sec:intro}

Post-training quantization (PTQ) of large-language-model planners for autonomous driving is
attractive for edge deployment: models like OpenDriveVLA~\cite{opendrivevla} (Qwen2.5-0.5B
planner on top of UniAD~\cite{uniad} perception) can in principle run on vehicle-grade compute at
a fraction of the FP16 memory footprint. Standard evaluation reports W4 as lossless---identical
open-loop nuScenes L2 to FP16.

We argue this conclusion is an evaluation artifact, not a safety guarantee, and provide evidence
that \emph{closed-loop} evaluation is required to surface quantization damage. Our contributions:
\begin{enumerate}
  \item An \textbf{ego-ablation} proving that nuScenes open-loop L2 is $\sim$entirely ego-motion
        extrapolation, making it insensitive to quantization damage on the perception/planning
        pathway (Sec.~\ref{sec:openloop}).
  \item A \textbf{closed-loop result on NeuroNCAP} (16 scenario instances, 10 seeds each, real
        neural-rendered cameras) showing that naive per-channel INT4 collapses generation
        coherence to 3.7\% while group-128 INT4 matches FP16 at 73.5\%
        (Sec.~\ref{sec:closedloop}).
  \item A \textbf{PTQ technique sweep} (per-channel, group-128, asymmetric, AWQ) confirming
        group-wise quantization as the minimal sufficient fix for the generation-collapse failure
        mode (Sec.~\ref{sec:sweep}).
  \item A \textbf{metric recommendation}: generation coherence (malformed trajectory rate) rather
        than collision rate, which is confounded by scene geometry and the chosen degenerate
        fallback (Sec.~\ref{sec:metric}).
\end{enumerate}

\section{Related Work}
\label{sec:related}

\textbf{VLA quantization.}
QuaRot~\cite{quarot} and SpinQuant~\cite{spinquant} enable low-bit activation quantization for
LLMs via orthogonal rotation. $\Omega$-QVLA~\cite{omegaqvla} and DyQ-VLA~\cite{dyqvla} extend
this to diffusion-head VLAs for manipulation, reporting near-lossless W4A4 on task-success
metrics. Group quantization (one scale per $g$ weights) is standard in
GPTQ~\cite{gptq} and AWQ~\cite{awq}; its safety implications in closed-loop driving are unstudied.

\textbf{Open-loop nuScenes evaluation.}
UniAD~\cite{uniad}, ST-P3~\cite{stp3}, and OpenDriveVLA~\cite{opendrivevla} all report open-loop
L2 on nuScenes. Prior work has noted that ego-motion extrapolation inflates open-loop scores
relative to closed-loop performance~\cite{closed_loop_critique}; we quantify this effect precisely
for the quantization setting.

\textbf{NeuroNCAP and closed-loop safety.}
NeuroNCAP~\cite{neuroncap} synthesizes safety-critical scenarios (frontal, side, and stationary
obstacle approaches) by substituting actors into real nuScenes scenes with NeuRAD neural
rendering~\cite{neurad}. The NCAP score (0--5; 5 = no collision) and the per-scenario collision
rate provide closed-loop safety metrics absent from open-loop benchmarks.

\section{Setup}
\label{sec:setup}

\textbf{Model.} OpenDriveVLA-0.5B~\cite{opendrivevla}: a Qwen2.5-0.5B planner conditioned on
UniAD stage-1 track+map perception, ego-state, and multi-camera image features.

\textbf{PTQ.} Weight-only fake-quant applied to all 168 Linear layers in
\texttt{model.model.layers} (357.8M / 500M parameters). Perception (UniAD FP32 ResNet),
projectors, embeddings, and lm\_head stay FP. Techniques: symmetric RTN per-output-channel;
RTN group-wise (group size 128 and 64); asymmetric RTN group-wise; AWQ~\cite{awq}
(calibrated on live rendered rollout activations). The model node snapshots FP weights once and
re-quantizes in-place via a \texttt{/quantize} REST endpoint for fast sweeps.

\textbf{Open-loop.} ST-P3~\cite{stp3} and UniAD-style L2 + collision on the 162 nuScenes-mini
$\cap$ val samples.

\textbf{Closed-loop.} Three-node NeuroNCAP harness: neuro-ncap orchestrator $+$ NeuRAD
renderer~\cite{neurad} (:8000) $+$ OpenDriveVLA model node (:9000). Scene selection: 14 scenes
(16 scenario instances: frontal, side, and stationary obstacles). Per-configuration: 10 paired
seeds with identical per-seed jitter across configs; only the quantization differs. Metric:
\emph{generation coherence} = fraction of predicted frames with a parseable $(x,y)$ trajectory
token (well-formed output rate). All-zero ``stay-in-place'' responses count as malformed
(degenerate).

\section{Open-Loop Is Blind}
\label{sec:openloop}

\subsection{W4 Equals FP16 on Open-Loop L2}

Table~\ref{tab:openloop} shows the nuScenes-mini open-loop results across weight precisions.
W8 and W4 are lossless (0.33\,m L2); W3 shows a $6\times$ L2 jump and 8\% malformed outputs;
W2 is catastrophic.

\begin{table*}[t]\centering
\caption{Open-loop planning on nuScenes-mini $\cap$ val (162 samples). ST-P3 L2 avg / malformed
output rate. W4 equals FP16; only W3 and W2 show degradation.}
\label{tab:openloop}
\begin{tabular}{lccc}
\toprule
Config & ST-P3 L2 (avg) & UniAD L2 (avg) & Malformed \\
\midrule
FP16 & 0.33\,m & 0.66\,m & 0\% \\
W8 & 0.33\,m & 0.66\,m & 0\% \\
W4 per-channel & 0.33\,m & 0.67\,m & 0\% \\
W3 per-channel & 1.91\,m & 3.26\,m & 8\% \\
W2 per-channel & 6.45\,m & $\sim$10\,m & 100\% \\
\bottomrule
\end{tabular}
\end{table*}

\subsection{Ego-Ablation: Why Open-Loop Cannot See Quantization Damage}
\label{sec:ablation}

We zero the ego-state (\texttt{gt\_ego\_lcf\_feat}) and history (\texttt{gt\_ego\_his\_trajs})
fed to the planner text prompt, while keeping all perception features:

\begin{table*}[t]\centering
\caption{Ego-state ablation on FP16 model (nuScenes-mini, 162 samples). Removing ego-state
collapses L2 $19\times$ while outputs remain well-formed.}
\label{tab:ego}
\begin{tabular}{lcc}
\toprule
Condition & ST-P3 L2 avg & Malformed \\
\midrule
Normal (ego + history + perception) & 0.33\,m & 0\% \\
Ego-state \& history zeroed (perception only) & 6.38\,m & 0\% \\
\bottomrule
\end{tabular}
\end{table*}

The $19\times$ L2 jump with a still-well-formed model reveals that nuScenes open-loop planning
is \emph{primarily ego-motion extrapolation from the velocity/history text}. The perception and
VLA pathway contribute little to this metric. Consequently, W4 quantization---which does not
damage the token-level ego-text processing---scores identically to FP16 on open-loop L2, even
if it harms the perception/planning pathway. Only closed-loop evaluation, where the model must
actually drive and perception is on the critical path, can surface this damage.

\section{Closed-Loop Generation Coherence}
\label{sec:closedloop}

\subsection{Per-Scene Results: 16 Scenario Instances, 10 Seeds}

Table~\ref{tab:perscene} reports the generation coherence rate per scenario instance for FP16,
naive per-channel W4, and group-128 W4.

\begin{table*}[t]\centering\small
\caption{Generation coherence (\%) per scenario instance (10 seeds each). FP16 vs.\ naive W4
(per-channel) vs.\ W4 group-128. Mean across all 16 instances (two no-collision scenarios
excluded from coherence mean): FP16 = 76.0\%, naive = 3.7\%, group-128 = 73.5\%.}
\label{tab:perscene}
\begin{tabular}{lccc}
\toprule
Scenario & FP16 & W4 naive & W4 g128 \\
\midrule
0099\_stationary & 60.6 & 0.0 & 30.0 \\
0101\_stationary & 75.0 & 14.3 & 46.7 \\
0103\_frontal & 74.3 & 0.0 & 72.9 \\
0106\_frontal & 66.9 & 0.0 & 78.7 \\
0108\_side & 85.7 & 0.0 & 91.6 \\
0108\_stationary & 72.4 & 0.0 & 74.7 \\
0110\_frontal & 70.7 & 0.0 & 74.3 \\
0278\_side & 75.3 & 0.0 & 67.9 \\
0278\_stationary & 79.7 & 0.0 & 83.5 \\
0331\_stationary & 83.2 & 5.0 & 97.5 \\
0346\_frontal & 83.1 & 15.6 & 86.9 \\
0783\_stationary & 72.4 & 0.0 & 73.1 \\
0796\_stationary & 72.4 & 8.6 & 57.9 \\
0921\_side & 90.0 & 14.2 & 90.4 \\
0923\_frontal & 74.6 & 0.8 & 72.1 \\
0966\_stationary & 79.4 & 0.0 & 78.3 \\
\midrule
\textbf{Mean} & \textbf{76.0} & \textbf{3.7} & \textbf{73.5} \\
\bottomrule
\end{tabular}
\end{table*}

Naive per-channel W4 collapses on 15 of 16 scenarios ($\leq$15\% coherence). Group-128 W4
recovers $\geq$85\% of FP16 coherence on 13 of 16. Partial failures on 3 scenarios (0099,
0101, 0796) are consistent with those scenes being hard even for FP16 ($<$75\%).

\subsection{Summary Table}

\begin{table*}[t]\centering
\caption{Generation coherence summary and technique sweep (well-formed predicted-frame rate,
scene-0103 frontal, 10 paired seeds).}
\label{tab:summary}
\begin{tabular}{lccc}
\toprule
Config & Coherence & Collapsed seeds & Notes \\
\midrule
FP16 & 71.4\% & 0/10 & baseline \\
\textbf{W4 per-channel (naive)} & \textbf{0.0\%} & \textbf{10/10} & total collapse \\
\textbf{W4 group-128} & \textbf{71.4\%} & \textbf{0/10} & $=$ FP16 \\
W4 group-64 & 71.4\% & 0/10 & $=$ FP16 \\
W4 asym group-128 & 71.4\% & 0/10 & $=$ FP16 \\
W3 per-channel & 21.4\% & 0/10 & partial \\
W3 group-128 & 85.7\% & 0/10 & $>$ FP16 (within noise) \\
W8 per-channel & 71.4\% & 0/10 & $=$ FP16 \\
W2 (all methods) & 0.0\% & 10/10 & precision floor \\
\bottomrule
\end{tabular}
\end{table*}

\section{PTQ Technique Sweep}
\label{sec:sweep}

Table~\ref{tab:summary} shows the full technique sweep. The key finding is that the collapse
is specific to \emph{per-channel} granularity: group-128, group-64, and asymmetric group-128 all
recover FP16 coherence at W4 precision. AWQ~\cite{awq} at W4 also recovers (71.4\%); the W3-AWQ
collapse is a suspected implementation bug in the low-bit AWQ path and is not reported.

The advantage of group-wise quantization \emph{grows} as bits drop: at W3, group-128 recovers
to 85.7\% while per-channel is only 21.4\%. At W2, all methods collapse (precision floor). This
monotone behavior---group $\gg$ per-channel, advantage widening at lower bits---is consistent
with the quantization granularity mechanism: finer groups prevent the single worst-channel scale
from dominating the entire layer's quantization error.

\section{Metric Recommendation}
\label{sec:metric}

\textbf{Why not collision rate.}
The ``40\% collision rate for naive W4'' reading of the multi-seed data is a confounded proxy.
Inspection of \texttt{trajectories.json} shows that naive W4 emits malformed/empty trajectory
text on 14/14 frames in every seed (100\% collapse), hitting the all-zero ``stay-in-place''
degenerate fallback. The 4/10 collision rate then depends on which seed geometries result in a
non-moving ego being struck, not on any planning precision effect. Two indicators:
\begin{enumerate}
  \item \textbf{Non-monotone:} W3-naive does NOT collapse (21.4\% coherence) but W4-naive does;
        lower bits cannot be ``safer'' under a genuine precision effect.
  \item \textbf{OOD-triggered:} On real nuScenes images (open-loop) naive-W4 is lossless; it
        collapses only on NeuRAD-rendered inputs (out-of-distribution, $\sim$0 detection recall).
\end{enumerate}

\textbf{Recommended metric.}
Report generation coherence (well-formed trajectory output rate) as the primary metric. It
directly measures whether the quantized model produces usable plans, is not mediated by the
degenerate fallback choice, and is computable from logged trajectory text without re-running
the simulation.

\section{Discussion and Limitations}
\label{sec:discussion}

\textbf{Mechanism.}
The generation collapse under per-channel INT4 is a \emph{quantization $\times$ domain-shift}
interaction. On real nuScenes images (in-distribution), per-channel W4 is well-formed; on
NeuRAD-rendered frames (out-of-distribution), the same quantization causes collapse. The
domain shift is not an artifact---deployed models will face distribution shift---but it means
the result strictly characterizes robustness to OOD inputs, not in-distribution degradation.

\textbf{Single scene limitation.}
The technique-sweep and multi-seed analysis use scene-0103 (the only available NeuRAD checkpoint).
The 14-scene coherence result extends to all released NeuroNCAP scenes but with only one seed
per scene for FP16 vs.\ per-channel vs.\ group-128; the 10-seed paired comparison is scene-0103
only. Generalization to additional scenes and a stronger renderer (requiring the gated 330\,GB
trainval dataset) is left to future work.

\textbf{Open-loop is sufficient for weight-only W3/W2 detection.}
The W3 and W2 degradations (6$\times$ and $>$20$\times$ L2 jump with 8\%/100\% malformed outputs)
are visible in open-loop L2. The blind spot is specifically the W4 per-channel vs.\ W4 group
distinction, which has no open-loop signature but a large closed-loop coherence effect.

\section{Conclusion}
\label{sec:conclusion}

Standard nuScenes open-loop L2 calls W4 quantization lossless because the metric is
$\sim$entirely ego-motion extrapolation (19$\times$ L2 jump on ego-ablation, well-formed outputs
throughout). Closed-loop NeuroNCAP evaluation across 16 real-camera scenario instances reveals
that naive per-channel W4 collapses generation coherence to 3.7\% (vs.\ 76.0\% for FP16) while
group-128 W4 restores FP16-level robustness (73.5\%). The fix is one parameter change;
group-64, asymmetric, and AWQ all recover equivalently. W8 and W4-group128 are deployment-safe;
naive W4 per-channel is not. Open-loop L2 cannot distinguish them.

\begin{thebibliography}{99}

\bibitem{opendrivevla}
  OpenDriveVLA Team, ``OpenDriveVLA: Driving with Vision-Language-Action Models,'' 2025.

\bibitem{uniad}
  Y.~Hu et al., ``Planning-Oriented Autonomous Driving,'' in \textit{CVPR}, 2023.

\bibitem{stp3}
  H.~Hu et al., ``ST-P3: End-to-End Vision-Based Autonomous Driving via Spatial-Temporal
  Feature Learning,'' in \textit{ECCV}, 2022.

\bibitem{neuroncap}
  W.~Ljungbergh et al., ``NeuroNCAP: Photorealistic Closed-Loop Safety Testing for Autonomous
  Driving,'' \textit{arXiv:2404.07762}, 2024.

\bibitem{neurad}
  A.~Tonderski et al., ``NeuRAD: Neural Rendering for Autonomous Driving,''
  in \textit{CVPR}, 2024.

\bibitem{quarot}
  S.~Ashkboos et al., ``QuaRot: Outlier-Free 4-Bit Inference in Rotated LLMs,'' 2024.

\bibitem{spinquant}
  Z.~Liu et al., ``SpinQuant: LLM Quantization with Learned Rotations,'' 2024.

\bibitem{gptq}
  E.~Frantar et al., ``GPTQ: Accurate Post-Training Quantization for GPT,'' 2022.

\bibitem{awq}
  J.~Lin et al., ``AWQ: Activation-aware Weight Quantization for LLM Compression,'' 2023.

\bibitem{omegaqvla}
  ``$\Omega$-QVLA: Robust Quantization for VLAs via Composite Rotation and Per-step Scaling,''
  2026.

\bibitem{dyqvla}
  ``DyQ-VLA: Temporal-Dynamic-Aware Quantization for Embodied VLAs,''
  \textit{arXiv:2603.07904}, 2026.

\bibitem{closed_loop_critique}
  K.~Chitta et al., ``Is Open-Loop Evaluation Enough for End-to-End Autonomous Driving?''
  in \textit{CVPR Workshops}, 2023.

\end{thebibliography}

\end{document}
